% =====================================================================
%  Gemeinsame Farben und Stile der Zeichnungen (TikZ, selbst erstellt)
%  In der Präambel einbinden: % =====================================================================
%  Gemeinsame Farben und Stile der Zeichnungen (TikZ, selbst erstellt)
%  In der Präambel einbinden: % =====================================================================
%  Gemeinsame Farben und Stile der Zeichnungen (TikZ, selbst erstellt)
%  In der Präambel einbinden: % =====================================================================
%  Gemeinsame Farben und Stile der Zeichnungen (TikZ, selbst erstellt)
%  In der Präambel einbinden: \input{zeichnungen/stil}
%  Lizenz wie das Dokument: CC BY-SA 3.0
% =====================================================================
\usetikzlibrary{calc,arrows.meta}

\definecolor{okgruen}{RGB}{46,160,67}
\definecolor{nokrot}{RGB}{207,34,46}
\definecolor{feld}{RGB}{244,246,248}
\definecolor{holz}{RGB}{226,196,150}
\definecolor{holzrand}{RGB}{150,110,60}
\definecolor{nut}{RGB}{186,150,100}
\definecolor{bediengruen}{RGB}{130,200,60}  % Grün der Bedienelemente

\tikzset{
  % Bauteile der Maschine
  gehaeuse/.style={fill=black!12, draw=black!75, line width=0.7pt, line join=round},
  griff/.style={fill=black!65, draw=black!85, line width=0.7pt, line join=round},
  metall/.style={fill=black!30, draw=black!75, line width=0.6pt, line join=round},
  werkstueck/.style={fill=holz, draw=holzrand, line width=0.6pt, line join=round},
  fraeser/.style={fill=black!55, draw=black!85, line width=0.5pt, line join=round},
  % Pfeile, Maße, Beschriftung
  vorschub/.style={-{Stealth[length=2.6mm, width=2.2mm]}, line width=1.3pt},
  drehung/.style={-{Stealth[length=1.8mm, width=1.6mm]}, line width=0.9pt},
  mass/.style={{Stealth[length=1.4mm]}-{Stealth[length=1.4mm]}, line width=0.4pt, black!80},
  hilfslinie/.style={line width=0.3pt, black!60},
  marke/.style={circle, draw=black, fill=white, line width=0.8pt,
    minimum size=4.4mm, inner sep=0pt, font=\scriptsize\bfseries},
  leitung/.style={line width=0.5pt, black!80},
  % Bewertung: \pic at (x,y) {ok}; bzw. {nok}
  ok/.pic={\fill[okgruen] (0,0) circle (3);
    \draw[white, line width=1.3pt, line cap=round, line join=round]
      (-1.4,0.1) -- (-0.4,-1.1) -- (1.5,1.2);},
  nok/.pic={\fill[nokrot] (0,0) circle (3);
    \draw[white, line width=1.3pt, line cap=round]
      (-1.1,-1.1) -- (1.1,1.1) (-1.1,1.1) -- (1.1,-1.1);},
}

% Feld mit hellem Hintergrund und Überschrift: \feld{x}{y}{Breite}{Höhe}{Titel}
\newcommand{\feld}[5]{%
  \fill[feld, rounded corners=2mm] (#1,#2) rectangle ++(#3,#4);
  \node[anchor=north west, font=\small\bfseries] at (#1+3,#2+#4-2) {#5};}

%  Lizenz wie das Dokument: CC BY-SA 3.0
% =====================================================================
\usetikzlibrary{calc,arrows.meta}

\definecolor{okgruen}{RGB}{46,160,67}
\definecolor{nokrot}{RGB}{207,34,46}
\definecolor{feld}{RGB}{244,246,248}
\definecolor{holz}{RGB}{226,196,150}
\definecolor{holzrand}{RGB}{150,110,60}
\definecolor{nut}{RGB}{186,150,100}
\definecolor{bediengruen}{RGB}{130,200,60}  % Grün der Bedienelemente

\tikzset{
  % Bauteile der Maschine
  gehaeuse/.style={fill=black!12, draw=black!75, line width=0.7pt, line join=round},
  griff/.style={fill=black!65, draw=black!85, line width=0.7pt, line join=round},
  metall/.style={fill=black!30, draw=black!75, line width=0.6pt, line join=round},
  werkstueck/.style={fill=holz, draw=holzrand, line width=0.6pt, line join=round},
  fraeser/.style={fill=black!55, draw=black!85, line width=0.5pt, line join=round},
  % Pfeile, Maße, Beschriftung
  vorschub/.style={-{Stealth[length=2.6mm, width=2.2mm]}, line width=1.3pt},
  drehung/.style={-{Stealth[length=1.8mm, width=1.6mm]}, line width=0.9pt},
  mass/.style={{Stealth[length=1.4mm]}-{Stealth[length=1.4mm]}, line width=0.4pt, black!80},
  hilfslinie/.style={line width=0.3pt, black!60},
  marke/.style={circle, draw=black, fill=white, line width=0.8pt,
    minimum size=4.4mm, inner sep=0pt, font=\scriptsize\bfseries},
  leitung/.style={line width=0.5pt, black!80},
  % Bewertung: \pic at (x,y) {ok}; bzw. {nok}
  ok/.pic={\fill[okgruen] (0,0) circle (3);
    \draw[white, line width=1.3pt, line cap=round, line join=round]
      (-1.4,0.1) -- (-0.4,-1.1) -- (1.5,1.2);},
  nok/.pic={\fill[nokrot] (0,0) circle (3);
    \draw[white, line width=1.3pt, line cap=round]
      (-1.1,-1.1) -- (1.1,1.1) (-1.1,1.1) -- (1.1,-1.1);},
}

% Feld mit hellem Hintergrund und Überschrift: \feld{x}{y}{Breite}{Höhe}{Titel}
\newcommand{\feld}[5]{%
  \fill[feld, rounded corners=2mm] (#1,#2) rectangle ++(#3,#4);
  \node[anchor=north west, font=\small\bfseries] at (#1+3,#2+#4-2) {#5};}

%  Lizenz wie das Dokument: CC BY-SA 3.0
% =====================================================================
\usetikzlibrary{calc,arrows.meta}

\definecolor{okgruen}{RGB}{46,160,67}
\definecolor{nokrot}{RGB}{207,34,46}
\definecolor{feld}{RGB}{244,246,248}
\definecolor{holz}{RGB}{226,196,150}
\definecolor{holzrand}{RGB}{150,110,60}
\definecolor{nut}{RGB}{186,150,100}
\definecolor{bediengruen}{RGB}{130,200,60}  % Grün der Bedienelemente

\tikzset{
  % Bauteile der Maschine
  gehaeuse/.style={fill=black!12, draw=black!75, line width=0.7pt, line join=round},
  griff/.style={fill=black!65, draw=black!85, line width=0.7pt, line join=round},
  metall/.style={fill=black!30, draw=black!75, line width=0.6pt, line join=round},
  werkstueck/.style={fill=holz, draw=holzrand, line width=0.6pt, line join=round},
  fraeser/.style={fill=black!55, draw=black!85, line width=0.5pt, line join=round},
  % Pfeile, Maße, Beschriftung
  vorschub/.style={-{Stealth[length=2.6mm, width=2.2mm]}, line width=1.3pt},
  drehung/.style={-{Stealth[length=1.8mm, width=1.6mm]}, line width=0.9pt},
  mass/.style={{Stealth[length=1.4mm]}-{Stealth[length=1.4mm]}, line width=0.4pt, black!80},
  hilfslinie/.style={line width=0.3pt, black!60},
  marke/.style={circle, draw=black, fill=white, line width=0.8pt,
    minimum size=4.4mm, inner sep=0pt, font=\scriptsize\bfseries},
  leitung/.style={line width=0.5pt, black!80},
  % Bewertung: \pic at (x,y) {ok}; bzw. {nok}
  ok/.pic={\fill[okgruen] (0,0) circle (3);
    \draw[white, line width=1.3pt, line cap=round, line join=round]
      (-1.4,0.1) -- (-0.4,-1.1) -- (1.5,1.2);},
  nok/.pic={\fill[nokrot] (0,0) circle (3);
    \draw[white, line width=1.3pt, line cap=round]
      (-1.1,-1.1) -- (1.1,1.1) (-1.1,1.1) -- (1.1,-1.1);},
}

% Feld mit hellem Hintergrund und Überschrift: \feld{x}{y}{Breite}{Höhe}{Titel}
\newcommand{\feld}[5]{%
  \fill[feld, rounded corners=2mm] (#1,#2) rectangle ++(#3,#4);
  \node[anchor=north west, font=\small\bfseries] at (#1+3,#2+#4-2) {#5};}

%  Lizenz wie das Dokument: CC BY-SA 3.0
% =====================================================================
\usetikzlibrary{calc,arrows.meta}

\definecolor{okgruen}{RGB}{46,160,67}
\definecolor{nokrot}{RGB}{207,34,46}
\definecolor{feld}{RGB}{244,246,248}
\definecolor{holz}{RGB}{226,196,150}
\definecolor{holzrand}{RGB}{150,110,60}
\definecolor{nut}{RGB}{186,150,100}
\definecolor{bediengruen}{RGB}{130,200,60}  % Grün der Bedienelemente

\tikzset{
  % Bauteile der Maschine
  gehaeuse/.style={fill=black!12, draw=black!75, line width=0.7pt, line join=round},
  griff/.style={fill=black!65, draw=black!85, line width=0.7pt, line join=round},
  metall/.style={fill=black!30, draw=black!75, line width=0.6pt, line join=round},
  werkstueck/.style={fill=holz, draw=holzrand, line width=0.6pt, line join=round},
  fraeser/.style={fill=black!55, draw=black!85, line width=0.5pt, line join=round},
  % Pfeile, Maße, Beschriftung
  vorschub/.style={-{Stealth[length=2.6mm, width=2.2mm]}, line width=1.3pt},
  drehung/.style={-{Stealth[length=1.8mm, width=1.6mm]}, line width=0.9pt},
  mass/.style={{Stealth[length=1.4mm]}-{Stealth[length=1.4mm]}, line width=0.4pt, black!80},
  hilfslinie/.style={line width=0.3pt, black!60},
  marke/.style={circle, draw=black, fill=white, line width=0.8pt,
    minimum size=4.4mm, inner sep=0pt, font=\scriptsize\bfseries},
  leitung/.style={line width=0.5pt, black!80},
  % Bewertung: \pic at (x,y) {ok}; bzw. {nok}
  ok/.pic={\fill[okgruen] (0,0) circle (3);
    \draw[white, line width=1.3pt, line cap=round, line join=round]
      (-1.4,0.1) -- (-0.4,-1.1) -- (1.5,1.2);},
  nok/.pic={\fill[nokrot] (0,0) circle (3);
    \draw[white, line width=1.3pt, line cap=round]
      (-1.1,-1.1) -- (1.1,1.1) (-1.1,1.1) -- (1.1,-1.1);},
}

% Feld mit hellem Hintergrund und Überschrift: \feld{x}{y}{Breite}{Höhe}{Titel}
\newcommand{\feld}[5]{%
  \fill[feld, rounded corners=2mm] (#1,#2) rectangle ++(#3,#4);
  \node[anchor=north west, font=\small\bfseries] at (#1+3,#2+#4-2) {#5};}
